\paragraph{SerializableDictionary}
\texttt{SerializableDictionary<TKey,TValue>} provides dictionary operations such as \texttt{Add}, \texttt{TryAdd}, \texttt{TryGetValue}, \texttt{Remove}, \texttt{AddRange}, and \texttt{RemoveWhere}. Serialized entries retain Inspector order. Duplicate keys are handled through \texttt{DuplicateKeyPolicy}: keep the first value, keep the last value, or report the duplicate without silently choosing one.

\paragraph{SerializableHashSet}
\texttt{SerializableHashSet<T>} stores unique values and provides \texttt{Contains}, \texttt{Add}, \texttt{Remove}, bulk updates, and predicate-based removal. Serialized order is retained for Inspector display.

\paragraph{SerializableLookup}
\texttt{SerializableLookup<TKey,TValue>} maps one key to multiple values. Use \texttt{Add}, \texttt{GetValues}, or \texttt{TryGetValues} when a dictionary value would otherwise need to be a manually managed list.

All three types derive from \texttt{SerializableCollectionBase}. They expose validation results, change/rebuild events, read-only interfaces, and explicit \texttt{Rebuild()} behavior. Null entries, malformed entries, and duplicate keys are represented as validation issues with a severity and an explanatory message.

\begin{lstlisting}
using System;
using MLGWorks.Utils.Helpers.Collections;
using UnityEngine;

[Serializable]
public sealed class ItemTags : SerializableLookup<string, string> { }

public sealed class InventoryData : MonoBehaviour
{
    [SerializeField] private ItemTags tags = new ItemTags();

    private void Reset()
    {
        tags.Clear();
        tags.Add("Sword", "Weapon");
        tags.Add("Sword", "Melee");
        tags.Add("Potion", "Consumable");
    }
}
\end{lstlisting}

The repository also contains a matching Inspector example. The \texttt{Editor} side adds the default property drawer, extension registration, and editor-only copy, paste, duplicate, and in-memory preset operations.
